% SSSUP-FONTE-STORICA: fonti-locali/originali/2014-2020.zip
% Anno 2018, esercizio 3. Trasposizione dal sorgente LaTeX storico.
% Verificare editorialmente la trasposizione prima di considerarla definitiva.
\renewcommand{\Anno}{2018}
\renewcommand{\NumeroProblema}{3}

\renewcommand{\Testo}{
Un'azienda del centro Italia produce mattonelle quadrate di lato 9 cm al prezzo $k_1$ ciascuna o una loro variante di lato 9,5 cm al prezzo unitario  $k_2>k_1$. La pavimentazione di una stanza consiste nel ricoprire l'intero suolo della stanza, con mattonelle o loro ritagli. La pavimentazione deve rispettare le seguenti regole:
\begin{itemize}
\item se due mattonelle si toccano allora o si toccano per un angolo oppure hanno in comune un intero lato;
\item solo al bordo della stanza possono essere usati ritagli di mattonelle;
\item tra il muro ed una mattonella intera può essere usato al più un ritaglio;
\item una mattonella può essere tagliata più volte.
\end{itemize}
A causa del consumo della lama della taglierina, ogni taglio ha un prezzo $k_3$. Un architetto deve fare pavimentare una stanza quadrata di lato 6 m e deve scegliere se acquistare mattonelle standard da 9 cm di lato oppure la variante da 9,5 cm di lato.
\begin{enumerate}
\item[(a)] Calcolare il costo minimo per pavimentare la stanza con mattonelle standard  ed il costo minimo per pavimentarla con le mattonelle variate e stabilire sotto quali ipotesi su $k_1$, $k_2$, $k_3$ la prima soluzione è più conveniente;
\item[(b)] Se ad ogni taglio c'è una probabilità dell'$1\%$ di rompere la mattonella (e doverla buttare via) a quanto ammonta il costo atteso della pavimentazione con mattonelle standard?
\end{enumerate}
}

\renewcommand{\Risposta}{
% SSSUP-BOZZA-AI: Risposta; autore=ChatGPT; data=2026-11-12; stato=DA-VERIFICARE
% SSSUP-PROVENIENZA: ChatGPT (OpenAI), sessione ProblemiSSSUP del 2026-11-12; derivato dalla traccia canonica; revisione umana pendente.
\textbf{Risultati parziali, non un ottimo globale certificato.}
Nel modello di tagli rettilinei su un pezzo alla volta, senza spessore
di taglio e con riuso dei ritagli, siano $C_9$ e $C_{9{,}5}$ i costi minimi.
Le costruzioni e le stime dimostrate nella soluzione danno
\[
 4445k_1\le C_9\le
 \min\{4489k_1+134k_3,\;4467k_1+157k_3,\;4445k_1+181k_3\},
\]
\[
 3989k_2\le C_{9{,}5}\le3991k_2+128k_3.
\]
In particolare $4445$ \`e il numero minimo di mattonelle da $9$ cm
che occorre acquistare, ma ci\`o non determina da solo il costo minimo.

(b) Ponendo $p=0{,}99$, per la posa da un angolo descritta nella soluzione,
con tagli indipendenti e addebitati anche se falliscono, il costo atteso \`e
\[
 4356k_1+132\frac{k_1+k_3}{p}
 +\frac{k_1+(1+p)k_3}{p^2}.
\]
Questa \`e l'attesa di una strategia specificata, non il minimo atteso
fra tutte le strategie. Restano da determinare l'ottimo globale di costo,
la regione completa di convenienza e l'eventuale ottimo atteso.
}
\renewcommand{\Aiutino}{
% SSSUP-BOZZA-AI: Aiutino; autore=ChatGPT; data=2026-11-12; stato=DA-VERIFICARE
% SSSUP-PROVENIENZA: ChatGPT (OpenAI), sessione ProblemiSSSUP del 2026-11-12; derivato dalla traccia canonica; revisione umana pendente.
Lavorare in centimetri e confrontare $600=66\cdot9+6$ con
$600=63\cdot9{,}5+1{,}5$.
Non \`e obbligatorio mettere la prima mattonella intera in un angolo:
la fascia residua pu\`o essere divisa fra due pareti opposte.
Contare separatamente mattonelle acquistate e tagli effettuati,
riutilizzando i ritagli. Per la parte aleatoria, distinguere una
mattonella che richiede un solo taglio da una che ne richiede due o pi\`u.
}
\renewcommand{\SoluzioneUfficiale}{}
\renewcommand{\SoluzioneRedazionale}{
% SSSUP-BOZZA-AI: SoluzioneRedazionale; autore=ChatGPT; data=2026-11-12; stato=DA-VERIFICARE
% SSSUP-PROVENIENZA: ChatGPT (OpenAI), sessione ProblemiSSSUP del 2026-11-12; derivato dalla traccia canonica; revisione umana pendente.
\textbf{Portata della soluzione e convenzioni.}
Questa \`e una soluzione parziale: fornisce costruzioni esplicite,
limiti inferiori e costi attesi di strategie determinate, ma non
certifica il minimo globale richiesto dal punto (a).
Assumiamo $k_1,k_2>0$ e $k_3\ge0$. Ogni taglio \`e rettilineo,
attraversa un solo pezzo e ha spessore trascurabile; non si impilano
mattonelle per tagliarle insieme. I ritagli possono essere riutilizzati.
Le pose qui considerate hanno lati paralleli alle pareti e sono
reticoli rettangolari, senza giunzioni a T: due pezzi adiacenti condividono
un intero lato. Tutti i pezzi tagliati sono sul bordo della stanza.
Non si assume, senza dimostrarlo, che fra queste pose vi sia l'ottimo
fra tutte quelle ammesse dal testo.

\textbf{(a) Limiti inferiori per area.}
La stanza misura $600\times600$ cm. Qualunque posa richiede almeno
\[
 \left\lceil\frac{600^2}{9^2}\right\rceil=4445
 \quad\hbox{oppure}\quad
 \left\lceil\frac{600^2}{9{,}5^2}\right\rceil=3989
\]
mattonelle acquistate. Poich\'e i tagli non hanno costo negativo,
\[
 C_9\ge4445k_1,\qquad C_{9{,}5}\ge3989k_2.
\]

\textbf{Mattonelle da $9$ cm: tre pose.}
Poich\'e $600=66\cdot9+6$, nelle tre costruzioni il blocco di mattonelle
intere \`e sempre formato da $66^2=4356$ mattonelle.

\emph{Posa A: partenza da un angolo.}
Il blocco intero occupa un quadrato $594\times594$ aderente a due pareti.
Restano $132$ rettangoli $9\times6$ e un quadrato $6\times6$.
Ciascun rettangolo si ottiene da una nuova mattonella con un taglio;
il quadrato richiede una nuova mattonella e due tagli. Si ha quindi
\[
 C_A=4489k_1+134k_3.
\]
Questi conteggi sono minimi per questa specifica posa: due rettangoli
$9\times6$, oppure uno di essi insieme al quadrato $6\times6$,
non possono provenire dalla stessa mattonella, gi\`a per ragioni di area.
Ogni rettangolo richiede almeno un taglio, il quadrato almeno due.

\emph{Posa B: fasce da $3$ cm su due pareti opposte.}
Centriamo il blocco intero in una direzione e lo appoggiamo a una parete
nell'altra. Servono $132$ rettangoli $9\times3$, $66$ rettangoli
$9\times6$ e due rettangoli $3\times6$ agli angoli.
I primi si ricavano da $44$ mattonelle, ciascuna divisa in tre strisce
con due tagli. I secondi richiedono $66$ mattonelle e $66$ tagli.
I due pezzi d'angolo si ricavano da una sola mattonella con tre tagli:
si estrae un rettangolo $6\times9$ e lo si divide in tre rettangoli
$6\times3$, utilizzandone due. Ne segue
\[
 C_B=4467k_1+157k_3.
\]

\emph{Posa C: fasce da $3$ cm su tutte le pareti.}
Centriamo il blocco $594\times594$ nella stanza. Servono $264$
rettangoli $9\times3$ e quattro quadrati $3\times3$.
I rettangoli provengono da $88$ mattonelle con $176$ tagli.
Per i quattro quadrati basta un'altra mattonella: due tagli estraggono
un quadrato $6\times6$ e altri tre lo dividono nei quattro quadrati
richiesti. Pertanto
\[
 C_C=4445k_1+181k_3.
\]
Questa posa raggiunge il limite inferiore di $4445$ mattonelle acquistate.
\`E dunque ottima per il solo numero di mattonelle, ma non per questo
\`e ottima per ogni costo dei tagli.
Ad esempio, per $k_1=k_3=1$ la posa A costa $4623$ e la posa C costa $4626$.

\textbf{Mattonelle da $9{,}5$ cm.}
Usiamo un blocco $63\times63$ di mattonelle intere, appoggiato a due
pareti, poich\'e $600=63\cdot9{,}5+1{,}5$.
Restano $126$ rettangoli $9{,}5\times1{,}5$ e un quadrato
$1{,}5\times1{,}5$.
Da ciascuna di $21$ mattonelle ricaviamo sei rettangoli con sei tagli,
lasciando una striscia di scarto larga $0{,}5$ cm.
Il quadrato finale richiede un'ulteriore mattonella e due tagli.
Abbiamo costruito una posa di costo
\[
 U_{9{,}5}=3991k_2+128k_3.
\]

\textbf{Confronti che si possono concludere rigorosamente.}
Poniamo $U_9=\min\{C_A,C_B,C_C\}$.
Se $U_9<3989k_2$, le mattonelle da $9$ cm sono certamente pi\`u
convenienti, perch\'e una loro posa costa meno di qualunque posa da
$9{,}5$ cm. Se invece $3991k_2+128k_3<4445k_1$, sono certamente
pi\`u convenienti le mattonelle da $9{,}5$ cm.
Fuori da queste regioni le stime non decidono il confronto.
Non sarebbe corretto presentare il solo confronto fra due costruzioni
come confronto fra i rispettivi ottimi globali.

\textbf{(b) Costo atteso di una strategia di taglio.}
Assumiamo che ciascun taglio riesca con probabilit\`a $p=0{,}99$,
indipendentemente dagli altri, e che anche i tagli falliti siano pagati.
Per ogni mattonella completiamo tutti i tagli prima di posarne i pezzi;
se un taglio fallisce, scartiamo tutti i pezzi di quella mattonella e
ripetiamo la stessa procedura su una nuova. Questo specifica anche
come si trattano i pezzi gi\`a ricavati prima di una rottura.

Se occorrono $m$ tagli riusciti, indichiamo con $A_m$ il costo atteso
per produrre il lotto di pezzi di una mattonella. Un tentativo costa
in media
\[
 k_1+k_3(1+p+\cdots+p^{m-1}),
\]
perch\'e il taglio numero $j$ viene eseguito soltanto se i $j-1$
precedenti sono riusciti. Il tentativo termina con successo con
probabilit\`a $p^m$; altrimenti si ricomincia. Dunque
\[
 A_m=k_1+k_3(1+p+\cdots+p^{m-1})+(1-p^m)A_m,
\]
ossia
\[
 \boxed{A_m=\frac{k_1+k_3(1+p+\cdots+p^{m-1})}{p^m}.}
\]
Le tre pose precedenti hanno quindi, con questa procedura, costi attesi
\[
 \begin{aligned}
 E_A&=4356k_1+132A_1+A_2,\\
 E_B&=4356k_1+66A_1+44A_2+A_3,\\
 E_C&=4356k_1+88A_2+A_5.
 \end{aligned}
\]
In particolare
\[
 E_A=4356k_1+132\frac{k_1+k_3}{p}
              +\frac{k_1+(1+p)k_3}{p^2}.
\]
Per $p=1$ si ritrovano esattamente i costi deterministici delle tre pose.
Il minimo di $E_A,E_B,E_C$ \`e il migliore fra questi tre procedimenti,
non il minimo atteso globale. Una soluzione completa richiederebbe
anche la minimizzazione sulle altre pose e sulle altre possibili
strategie di taglio e riutilizzo; qui tale ottimalit\`a non \`e dimostrata.
}
\renewcommand{\SoluzioniAlternative}{}
\renewcommand{\Osservazioni}{}
\renewcommand{\Commenti}{}
